%\section*{Contents}

Chapter I.  \textbf{Introduction. Articles 1-27. page 1}
Sound due to Vibrations. Finite velocity of Propagation. Velocity independent
of Pitch. Regnault’s experiments. Sound propagated in water.  Wheatstone’s
experiment. Enfeeblement of Sound by distance. Notes and Noises. Musical notes
due to periodic vibrations. Siren of Cagniard de la Tour.  Pitch dependent upon
Period. Relationship between musical notes. The same ratio of periods
corresponds to the same interval in all parts of the scale. Harmonic scales.
Diatonic scales. Absolute Pitch. Necessity of Temperament. Equal Temperament.
Table of Frequencies. Analysis of Notes. Notes and Tones. Quality dependent
upon harmonic overtones. Resolution of Notes by ear uncertain. Simple tones
correspond to simple pendulous vibrations.

Chapter II.  \textbf{Harmonic Motions.  Articles 28-42a. page 18}
Composition of harmonic motions of like period. Harmonic Curve. Composition of
two vibrations of nearly equal period. Beats. Fourier’s Theorem.
[32a Beats of approximate consonances.]
Vibrations in perpendicular directions. Lissajous’ Cylinder.  Lissajous’ Figures.
Blackburn’s pendulum. Kaleidophone. Optical methods of composition and
analysis. The vibration microscope. Intermittent Illumination.
{[} Resultant of a large number of vibrations whose phases are accidentally distributed.{]}

Chapter III.  \textbf{Systems Having One Degree of Freedom.  Articles 43-68d page 35}
Systems with one degree of freedom. Independence of amplitude and period.
Frictional force proportional to velocity. Forced vibrations.  Principle of
Superposition. Beats due to superposition of forced and free vibrations.
Various degrees of damping. String with Load. Method of Dimensions. Ideal Tuning
fork. Forks give nearly pure tones.  Forks as standards of pitch.
[54a Dependence upon temperature.]
[54b Slow versus quick beats.]
Scheibler’s methods of tuning. Scheibler’s Tonometers. Compound Pendulum. Forks
driven by electromagnetism.
[58a  Phonic wheel.]
Fork Interrupter. Resonance.
[60a Intermittent vibrations.]
General solution for one degree of freedom.
[61a Instability.]
61b Terms of the second order give rise to derived tones.
[68c  Maintenance]
[68d Methods of determining absolute pitch.]


Chapter IV.  \textbf{Vibrating Systems in General.  Articles 69-95 page 67}
Generalized coordinates. Expression for potential energy. Statical theorems.
Initial motions. Expression for kinetic energy. Reciprocal
theorem. Thomson’s
[-1 Kelvin's]
theorem. Lagrange’s equations. The dissipation
function. Coexistence of small motions. Free vibrations without friction.
Normal coordinates. The free periods fulfil a stationary condition. An
accession of inertia increases the free periods. A relaxation
of spring increases the free periods. The greatest free period is an
absolute maximum. Hypothetical types of vibration. Example from
string. Approximately simple systems. String of variable density.
Normal functions. Conjugate property.
[Introduction of one constraint.]
[Several constraints.]
Determination of constants to suit arbitrary initial conditions. Stokes’ theorem.


Chapter V.
\textbf{Vibrating Systems in General Continued.  Articles 96-117 page 97}
Cases in which the three functions T. F. V are simultaneously reducible to
sums of squares. Generalization of Young’s theorem on the nodal
points of strings. Equilibrium theory. Systems started from rest as
deflected by a force applied at one point. Systems started from the
equilibrium configuration by an impulse applied at one point. Systems
started from rest as deflected by a force uniformly distributed. Influence of
small frictional forces on the vibrations of a system. Solution of
the general equations for free vibrations.
[Routh’s theorems.]
[Instability.]
Impressed Forces. Principle
of the persistence of periods. Inexorable motions. Reciprocal Theorem.
Application to free vibrations. Statement of reciprocal theorem
for harmonic forces. Applications. Extension to cases in which the
constitution of the system is a function of the period.
[Reaction at driving point.]
Equations for two degrees of freedom. Roots of determinantal equation. Intermittent
vibrations. March of periods as inertia is gradually increased.
Reaction of a dependent system.

Chapter VI. \textbf{Articles 118—148c page 127}
Law of extension of a string. Transverse vibrations. Solution of the problem
for a string whose mass is concentrated in equidistant points. Derivation of
solution for continuous string. Partial differential equation.  Expressions for
V and T. Most general form of simple harmonic motion. Strings with fixed
extremities. General motion of a string periodic. Mersenne’s Laws.  Sonometer.
Normal modes of vibration.  Determination of constants to suit arbitrary
initial circumstances. Case of plucked string. Expressions for T and V in terms
of normal coordinates.  Normal equations of motion. String excited by
plucking.  Young’s theorem. String excited by an impulse. Problem of pianoforte
string.  Friction proportional to velocity. Comparison with equilibrium theory.
Periodic force applied at one point. Modifications due to yielding of the
extremities. Proof of Fourier’s theorem. Effects of a finite load. Correction
for rigidity. Problem of violin string.  Strings stretched on curved surfaces.
Solution for the case of the sphere. Correction for irregularities of density.
[ -1 Arbitrary displacementof every period.]
Theorems of Sturm and Liouville for a string of variable density.
[ -1 Density proportional \(x^{-2}\)]
[ -1 Nodes of forced vibrations.]
Propagation
of waves along an unlimited string. Positive and negative waves. Stationary
Vibrations. Reflection at a fixed point. Deduction of solution for finite
string. Graphical method. Progressive wave with friction.
[148a Reflection at a junction of two strings.]
[148b Gradual transition.]
[148c Effect of imperfect flexibility.]

Chapter VII \textbf{Articles 149—159 page 188}
Classification of the vibrations of Bars. Differential equation for
longitudinal vibrations. Numerical values of the constants for steel. Solution
for a bar free at both ends. Deduction of solution for a bar with one end free.
and one fixed. Both ends fixed. Influence of small load.  Solution of problem
for bar with large load attached.
[-1 Reflection at a junction.]
Correction for lateral motion. Savart’s
\textit{son rauque}. Differential equation for torsional vibrations. Comparison
of velocities of longitudinal and torsional waves.

Chapter VIII. \textbf{Articles 160—192a page 201}
Potential energy of bending. Expression for kinetic energy. Derivation of
differential equation. Terminal conditions. General solution for a harmonic
vibration. Conjugate property of the normal functions.  Values of integrated
squares. Expression of V in terms of normal co ordinates. Normal equations of
motion. Determination of constants to suit initial conditions.
[-1 Numerical formulae for tuning-folks.]
Case of rod
started by a blow. Rod started from rest as deflected by a lateral force. In
certain cases the series of normal functions ceases to converge. Form of the
normal functions for a free-free bar. Laws of dependence of frequency on length
and thickness. Case when both ends are clamped. Normal functions for
clamped-free bar. Calculation of periods. Comparisons of pitch. Discussion of the
gravest mode of vibration of a free-free bar. Three nodes. Four nodes.  Gravest
mode for clamped-free bar. Position of nodes. Supported bar.  Calculation of
period for clamped-free bar from hypothetical type. Solution of problem for a bar
with a loaded end.  Effect of additions to a bar. Influence of irregularities of
density.  Correction for rotatory inertia. Roots of functions derived linearly
from normal functions. Formation of equation of motion when there is permanent
tension. Special terminal conditions. Resultant of two trains of waves of
nearly equal period. Fourier’s solution of problem for infinite bar.
[192a Circular Ring.]

Chapter IX. \textbf{Articles 193—213a page 250}
Tension of a membrane. Equation of motion. Fixed rectangular boundary.
Expression for V and T in terms of normal coordinates. Normal equations of
vibration. Examples of impressed forces. Frequency for an elongated rectangle
depends mainly on the shorter side. Cases in which different modes of vibration
have the same period. Derived modes thence arising. Effect of slight
irregularities. An irregularity may remove indeterminateness of normal modes.
Solutions applicable to a triangle. Expression of the general differential
equation by polar coordinates. Of the two functions. which occur in the
solution. one is excluded by the condition at the pole. Expressions for
Bessel’s equations. Formulae relating thereto. Table of the first two
functions.  Fixed circular boundary. Conjugate property of the normal functions
without restriction of boundary. Values of integrated squares. Expressions for
T and V in terms of normal functions. Normal equations of vibration for
circular membrane. Special case of free vibrations. Vibrations due to a
harmonic force uniformly distributed.
[-1 Force applied locally at the centre.]
Pitches of the various simple tones.
Table of the roots of Bessel’s functions. Nodal Figures. Circular membrane with
one radius fixed.  Bessel’s functions of fractional order. Effect of small
load. Vibrations of a membrane whose boundary is approximately circular. In
many cases the pitch of a membrane may be calculated from the area alone.  Of
all membranes of equal area that of circular form has the gravest pitch. Pitch
of a membrane whose boundary is an ellipse of small eccentricity. Method of
obtaining limits in cases that cannot be dealt with rigorously. Comparison of
frequencies in various cases of membranes of equal area. History of the
problem. Bourget’s experimental investigations.
[213a Kettle drums. Nodal curves of forced vibrations.]

Chapter X. \textbf{Articles 214-235a page 298}
Vibrations of Plates. Potential Energy of Bending. Transformation of 5V.
Superficial differential equation. Boundary conditions. Conjugate property of
normal functions. Transformation to polar coordinates. Form of general solution
continuous through pole. Equations determining the periods for a free
circular plate. Kirchhoff’s calculations.  Comparison with observation. Radii
of nodal circles. Irregularities give rise to beats. Generalization of
solution.
[Oscillation of nodes.]
Case of clamped. or supported. edge.
[Telephone plate.]
Disturbance of Chladni’s figures.
[Movements of sand.]
History of problem.  Mathieu’s criticisms. Rectangular plate with
supported edge. Rectangular plate with free edge. Boundary conditions. One
special case \(\mu = 0\) is amenable to mathematical treatment.  Investigation
of nodal figures. Wheatstone’s application of the method of superposition.
Comparison of Wheatstone’s figures with those really applicable to a plate in
the case ~4=0. Gravest mode of a square plate. Calculation of period on
hypothetical type. Nodal figures inferred from considerations of symmetry.
Hexagon. Comparison between circle and square.  Law connecting pitch and
thickness. In the case of a clamped edge any contraction of the boundary raises
the pitch. No gravest form for a free plate of given area. In similar plates
the period is as the linear dimension.  Wheatstone’s experiments on wooden
plates. Rudolph Koenig’s experiments. Vibrations of cylinder. or ring.  Motion
tangential as well as normal. Relation between tangential and normal motions.
Expressions for kinetic and potential energies. Equations of vibration.
Frequencies of tones. Comparison with Chladni.
[Fenknor’s observations.]
Tangential friction excites
tangential motion. Experimental verification. Beats due to irregularities.
[235a Examples of glass bells.  Church bells.]


CHAPTER Xa.
§235b -- 235 h CURVED PLATES OR SHELLS 395
Extensional Vibrations. Frequency independent of thickness. Inextensional
or flexural vibrations. Frequency proportional to thickness.
General conditions of inextension. Surface of second degree. Application to
sphere. Principal extensions of cylindrical surface. Potential energy.
Frequencies of extensional vibrations. Plane plate. Other
particular cases of cylinder. Potential energy of bending. Sphere.
Plane plate. Potential energy for cylindrical shell. Statical problems.
Frequency of flexural vibrations of cylindrical shell. Extensional
vibration of spherical shell. Flexural vibrations of spherical shell.
Normal modes. Potential energy. Kinetic energy. Frequencies in case
of hemisphere. Saucer of 120 degrees. References.

CHAPTER Xb.
§ 235 i — 235 y Electrical Vibrations. 438
Calculation of periods. Forced vibrations. Insertion of a leyden equivalent
to a negative inductance. Initial currents in a secondary circuit. Inversely
as the number of windings. Reaction of secondary circuit.
Train of circuits. Initial currents alternately opposite in sign.
Persistences. Resistance and inductance of two conductors in parallel.
Extreme values of frequency. Contiguous wires. Several conductors in
parallel.  Induction balance. Theory for simple harmonic currents.Two
conditions necessary for balance. Wheatstone’s bridge. Generalized
resistance. Current in bridge. Approximate balance. Hughes' arrangement.
Interrupters. Inductometers. Symmetrical arrangement.
Electromagnetic screening. Cylindrical conducting core. Time-constant of
free currents Induced electrical vibrations. Reaction upon
primary circuit. Induced currents in a wire. Maxwell's formulas.
Impedance. Kelvin’s theory of cables. Heaviside’s generalization.
Attenuation and distortion of signals. Bell’s telephone. Push and
pull theory. Experiment upon bipolar telephone. Minimum current
audible Microphone.


